% texts/books.tex -- THE BOOK REGISTRY, shared by every course on this machinery.
%
% Everything here answers one question: WHICH BOOKS CAN A DOCUMENT NAME, AND
% WHERE DO THEY LIVE? Nothing else belongs in this file.
%
% WHY IT IS IN THE SHARED MACHINERY. A book's title, publisher and URL are facts
% about the outside world: the same in every course that cites the book, and in
% every term. So they live here once, and a lecture in ANY course can point a
% student at algtrig without that course first having to declare it. A course
% that wants a different registry puts its own books.tex in its course-info/,
% which is searched first -- but then it gets ONLY its own, not both.
%
% WRITTEN FOR MATH 1003, and the comments below still speak from there: "this
% course" and the three documents named are MATH 1003's. It uses
% three books, which is more than usual:
%
%   apex      APEX Calculus            the main text -- limits onward
%   openstax  Calculus, Volume 1       actually used -- the functions review
%   algtrig   Algebra & Trigonometry   extra help, never taught from
%
% HOW IT IS LOADED, AND WHY THAT IS DIFFERENT FROM apex.tex. No master \inputs
% this file, and no master names it either: TextRef loads it unconditionally, in
% EVERY document, before the preamble ends. .latexmkrc puts the machinery on
% TEXINPUTS with a trailing //, so it is found BY NAME at any depth, the same
% way the text layers are.
%
% That "every document, whole" is the whole point of the file existing. Compare
% \usetext, which loads exactly ONE text layer, because vocabulary is exclusive:
% a document cannot call worked items "Example" and "Checkpoint" at once. A
% book's NAME and URL are not exclusive in the least, and three ordinary
% documents in this course prove it:
%
%   - the SYLLABUS names all three books, so it can be built on none of them
%     (it has no \usetext, and there is no course default to fall back on);
%   - ADVICE FOR STUDENTS points at algtrig throughout and at neither calculus
%     book -- and no document will ever say \usetext{algtrig}, because no
%     document IS an algtrig document;
%   - the FUNCTIONS LECTURES are built on OpenStax and still send students to
%     algtrig for the algebra underneath.
%
% Before this file existed, all three were hand-typed \href URLs with the domain
% baked in -- thirteen of them, each rotting on its own schedule, which is the
% exact failure TextRef.sty was written to prevent.
%
% THE MEMBERSHIP TEST, same as the text layer's: would you look at the book to
% decide the value? Its title, its publisher, where it lives, which page a
% section sits on -- yes. How titles are set, what a "Key Idea" is called, the
% noun for a worked item -- NO. The first is identity and lives here; the second
% is vocabulary and is exclusive, so it lives in apex.tex / openstax-calculus1.tex
% beside this file.
%
% WHICH IS WHY algtrig HAS NO FILE OF ITS OWN. It has identity and nothing else.


%%%%%%%%%%%%%%%%%%%%%%%%%%%%%%%%%%%%%%%%
% 1. THE MAIN TEXT
%%%%%%%%%%%%%%%%%%%%%%%%%%%%%%%%%%%%%%%%
% This is the STANDARD, all-in-one APEX web build. ULeth also publishes
% "accelerated" and split "standard" builds at sibling URLs whose chapters are
% ORDERED DIFFERENTLY, so their section numbers do not agree with the ones the
% lectures cite. Adopting one of those is a different base AND different pages.
%
% Keep the trailing slash on base -- slugs are appended verbatim.
%
% No publisher: APEX is Gregory Hartman's, published under a Creative Commons
% licence and rebuilt by the University of Lethbridge, and no one of those is
% the answer to "who publishes it". The syllabus says so in prose instead.
%
% No home: it defaults to base, which is right here -- the book's web edition is
% its front door.
\DeclareBook{apex}{%
  short = APEX,
  title = APEX Calculus,
  base  = https://opentext.uleth.ca/apex-calculus/,
}


%%%%%%%%%%%%%%%%%%%%%%%%%%%%%%%%%%%%%%%%
% 2. THE SECOND TEXT, FOR THE FUNCTIONS REVIEW
%%%%%%%%%%%%%%%%%%%%%%%%%%%%%%%%%%%%%%%%
% Volume 1 of the three-volume OpenStax Calculus. Note that its page slugs carry
% their own "pages/" segment and NO .html extension, unlike APEX's -- slug form
% is per publisher, so do not pattern-match one book's from the other's.
%
% home differs from base: OpenStax gives every book a catalogue page separate
% from the reading edition, and the catalogue page is what a student should be
% sent to first (it offers the PDF and print options as well).
\DeclareBook{openstax}{%
  short     = OpenStax Calculus,
  % Braced because of the comma: keyval would otherwise read "Volume 1" as a
  % key of its own and stop with an undefined-key error. (OpenStax's own
  % spelling drops the comma; the braces are here so the next title that needs
  % one has the idiom in front of it.)
  title     = {Calculus, Volume 1},
  publisher = OpenStax,
  home      = https://openstax.org/details/books/calculus-volume-1,
  base      = https://openstax.org/books/calculus-volume-1/,
}


%%%%%%%%%%%%%%%%%%%%%%%%%%%%%%%%%%%%%%%%
% 3. THE PREREQUISITE TEXT
%%%%%%%%%%%%%%%%%%%%%%%%%%%%%%%%%%%%%%%%
% Not taught from, ever. It is where a student is sent when the calculus is
% sound and the algebra underneath is not -- which is the commonest way to lose
% marks in this course, and the reason Advice for Students leans on it so hard.
%
% This is the book that could not exist under the old arrangement: identity with
% no vocabulary, so no \usetext file, so no base, so every reference to it was a
% full hand-typed URL.
\DeclareBook{algtrig}{%
  short     = Algebra \& Trigonometry,
  title     = Algebra and Trigonometry 2e,
  publisher = OpenStax,
  home      = https://openstax.org/details/books/algebra-and-trigonometry-2e,
  base      = https://openstax.org/books/algebra-and-trigonometry-2e/,
}


%%%%%%%%%%%%%%%%%%%%%%%%%%%%%%%%%%%%%%%%
% 4. SECTIONS OF algtrig THAT THIS COURSE POINTS AT
%%%%%%%%%%%%%%%%%%%%%%%%%%%%%%%%%%%%%%%%
% These are named-book \DeclareSection entries, and they live HERE rather than
% in a text layer for the reason §3 gives: no document is built on this book, so
% no \usetext would ever load a file holding them.
%
% Contrast the entries in apex.tex and openstax-calculus1.tex, which take no book name.
% The empty name means "the book THIS DOCUMENT is built on", which is the only
% thing a text layer can honestly say -- and the two namespaces deliberately do
% not cross. See TextRef.sty under \DeclareSection.
%
% DECLARE ONLY WHAT YOU CITE. Every line below is cited by a live document; the
% sparseness is a CHECK, not tidiness. \textbookref reports a key it cannot
% resolve and prints the reference unlinked, so mistyping 3.2 for 3.3 is
% something you SEE. Declare the whole book and that same typo silently emits a
% perfect-looking link to the wrong section. One entry per line, so these stay
% greppable without compiling anything.
%
% !! THE SLUGS ARE AS VERIFIED AS WHAT THEY REPLACED, AND NO MORE. Each was
% !! lifted verbatim from a working \href in advice-for-students.tex or a
% !! functions lecture, so this file introduced no new 404 risk -- but it
% !! inherited whatever was already wrong. LaTeX warns when a section KEY does
% !! not resolve; a slug that is simply wrong still produces a perfectly
% !! good-looking link to a 404. Check one by hand before leaning on it.
%
% TOPICS.tex's set of [xN.M] pointers is the authoritative list of INTENDED ones and
% is longer than this. Add a line here when a document actually cites it.

% --- Chapter 1: Prerequisites ---
\DeclareSection[algtrig]{1.2}{pages/1-2-exponents-and-scientific-notation}       % Exponents and Scientific Notation
\DeclareSection[algtrig]{1.3}{pages/1-3-radicals-and-rational-exponents}         % Radicals and Rational Exponents
\DeclareSection[algtrig]{1.6}{pages/1-6-rational-expressions}                    % Rational Expressions

% --- Chapter 3: Functions ---
\DeclareSection[algtrig]{3.1}{pages/3-1-functions-and-function-notation}         % Functions and Function Notation
\DeclareSection[algtrig]{3.2}{pages/3-2-domain-and-range}                        % Domain and Range
\DeclareSection[algtrig]{3.3}{pages/3-3-rates-of-change-and-behavior-of-graphs}  % Rates of Change and Behavior of Graphs
\DeclareSection[algtrig]{3.4}{pages/3-4-composition-of-functions}                % Composition of Functions
\DeclareSection[algtrig]{3.5}{pages/3-5-transformation-of-functions}             % Transformation of Functions

% --- Chapter 4: Linear Functions ---
\DeclareSection[algtrig]{4.1}{pages/4-1-linear-functions}                        % Linear Functions

% --- Chapter 5: Polynomial and Rational Functions ---
\DeclareSection[algtrig]{5.1}{pages/5-1-quadratic-functions}                     % Quadratic Functions
% 5.2 and 5.3 are the two entries in TOPICS.tex's OS 1.2 block that had no line here,
% so OS1.2's prepwork could not cite them. Unlike every slug above, these two
% were NOT lifted from a working \href -- they were fetched and checked against
% the live site (both 200, both titles matching) before being added.
\DeclareSection[algtrig]{5.2}{pages/5-2-power-functions-and-polynomial-functions} % Power Functions and Polynomial Functions
\DeclareSection[algtrig]{5.3}{pages/5-3-graphs-of-polynomial-functions}         % Graphs of Polynomial Functions
\DeclareSection[algtrig]{5.5}{pages/5-5-zeros-of-polynomial-functions}           % Zeros of Polynomial Functions
\DeclareSection[algtrig]{5.7}{pages/5-7-inverses-and-radical-functions}          % Inverses and Radical Functions

% --- Chapter 6: Exponential and Logarithmic Functions ---
\DeclareSection[algtrig]{6.1}{pages/6-1-exponential-functions}                   % Exponential Functions
\DeclareSection[algtrig]{6.3}{pages/6-3-logarithmic-functions}                   % Logarithmic Functions
\DeclareSection[algtrig]{6.5}{pages/6-5-logarithmic-properties}                  % Logarithmic Properties

% --- Chapter 7: The Unit Circle ---
\DeclareSection[algtrig]{7.1}{pages/7-1-angles}                                  % Angles
\DeclareSection[algtrig]{7.3}{pages/7-3-unit-circle-sine-and-cosine-functions}   % Unit Circle: Sine and Cosine Functions

% --- Chapter 8: Periodic Functions ---
\DeclareSection[algtrig]{8.1}{pages/8-1-graphs-of-the-sine-and-cosine-functions} % Graphs of the Sine and Cosine Functions
\DeclareSection[algtrig]{8.3}{pages/8-3-inverse-trigonometric-functions}         % Inverse Trigonometric Functions

% --- Chapter 9: Trigonometric Identities and Equations ---
\DeclareSection[algtrig]{9.1}{pages/9-1-solving-trigonometric-equations-with-identities} % Solving Trigonometric Equations with Identities
\DeclareSection[algtrig]{9.2}{pages/9-2-sum-and-difference-identities}           % Sum and Difference Identities
